%% ============================================================
%% Appendix: the order-4 symmetry of the dimension-4 verifier.
%% Include with \appendix \input{appendix-symmetry} in Dim4SQI.tex.
%% Uses \Echl and \Ecom from the main preamble.
%% ============================================================
\section{The order-4 symmetry of the verifier's chain}
\label{app:symmetry}

This appendix collects what can be proved about the automorphism of order
4 mentioned in Section~\ref{sec:conc}, and states precisely what remains
open before MIKE's saving of two kernel points out of four~\cite[Appendix
E.2]{BDD+26} could be applied to the verifier of Section~\ref{sec:dim4}.

\paragraph{Notation}
Throughout, $E$ denotes a supersingular elliptic curve over
$\mathbb{F}_{p^2}$ and $E^2 = E \times E$. Homomorphisms between squares of
curves are written as $2 \times 2$ matrices whose entries are homomorphisms
between the curves, acting on column vectors:
\[
\begin{pmatrix} u & v \\ w & z \end{pmatrix}
\begin{pmatrix} P \\ Q \end{pmatrix}
=
\begin{pmatrix} u(P) + v(Q) \\ w(P) + z(Q) \end{pmatrix}.
\]
An integer $n$ stands for the multiplication map $[n]$. For a homomorphism
$\psi$ we write $\widehat{\psi}$ for its dual. With respect to the product
polarisations, the dual of a matrix is the transpose of the matrix of
duals; in particular the dual of an integer matrix is its transpose. We
write $\Echl$ for the challenge curve, $\Ecom$ for the commitment curve,
$\sigma \colon \Echl \to \Ecom$ for the response, of odd degree $q$, and
$a_1, a_2$ for integers with $a_1^2 + a_2^2 = 2^e - q$, as in
Section~\ref{sec:dim4}. Finally
\[
A_1 = \Echl^2, \qquad A_2 = \Ecom^2, \qquad A = A_1 \times A_2,
\]
so that $A$ is the abelian fourfold on which the verifier's chain starts.

\subsection{The automorphism \texorpdfstring{$J$}{J} of \texorpdfstring{$E^2$}{E\^{}2}}

\begin{definition}
    On any square $E^2$, let
    \[
    J = \begin{pmatrix} 0 & 1 \\ -1 & 0 \end{pmatrix},
    \qquad\text{that is,}\qquad
    J(P, Q) = (Q, -P),
    \]
    and on $A = A_1 \times A_2$ let $D = \operatorname{diag}(J, J)$ be $J$
    acting on each square separately.
\end{definition}

\begin{lemma}
    \label{lem:J}
    \leavevmode
    \begin{enumerate}[label=(\roman*)]
        \item $J^2 = -1$, so $J$ has order 4.
        \item The endomorphism
        $\alpha = \left(\begin{smallmatrix} a_1 & a_2 \\ -a_2 & a_1
        \end{smallmatrix}\right)$ of Figure~\ref{fig:kani} equals
        $a_1 + a_2 J$, and its dual is
        $\widehat{\alpha} = \alpha^{\mathsf{T}} = a_1 - a_2 J$.
        \item $\widehat{\alpha}\,\alpha = \alpha\,\widehat{\alpha} = a_1^2 + a_2^2
        = 2^e - q$.
        \item $\widehat{J} = J^{\mathsf{T}} = -J = J^{-1}$, so $J$ preserves the
        product polarisation of $E^2$: it is an automorphism of $E^2$ as a
        principally polarised abelian surface.
    \end{enumerate}
\end{lemma}

\begin{proof}
    (i) $J^2(P, Q) = J(Q, -P) = (-P, -Q)$. (ii) We have
    \[
    (a_1 + a_2 J)(P, Q) = \bigl([a_1]P + [a_2]Q, \; -[a_2]P + [a_1]Q\bigr),
    \]
    which is $\alpha(P, Q)$; the dual
    of an integer matrix is its transpose, and $J^{\mathsf{T}} = -J$. (iii)
    $(a_1 + a_2 J)(a_1 - a_2 J) = a_1^2 - a_2^2 J^2 = a_1^2 + a_2^2$, since
    $J$ commutes with integers. (iv) $J^{\mathsf{T}} J = 1$ by direct
    computation, and an automorphism $\gamma$ preserves the polarisation
    exactly when $\widehat{\gamma}\gamma = 1$.
\end{proof}

\subsection{The verifier's isogeny commutes with \texorpdfstring{$D$}{D}}

Write $\Sigma = \operatorname{diag}(\sigma, \sigma) \colon A_1 \to A_2$, the
map $\sigma \oplus \sigma$ of Figure~\ref{fig:kani}, so that
$\widehat{\Sigma} = \operatorname{diag}(\widehat{\sigma},
\widehat{\sigma})$. Up to the order and signs of its blocks, which do not
affect anything below, the verifier's dimension-4 isogeny is the Kani
endomorphism~\cite[Section 4]{DLRW24}
\[
F = \begin{pmatrix} \alpha & \widehat{\Sigma} \\
    -\Sigma & \widehat{\alpha} \end{pmatrix}
\colon A_1 \times A_2 \longrightarrow A_1 \times A_2,
\qquad
F(x_1, x_2) = \bigl(\alpha(x_1) + \widehat{\Sigma}(x_2), \;
-\Sigma(x_1) + \widehat{\alpha}(x_2)\bigr),
\]
where the $\alpha$ in the top-left corner acts on $A_1$ and the
$\widehat{\alpha}$ in the bottom-right corner on $A_2$. Its dual is
\[
\widehat{F} = \begin{pmatrix} \widehat{\alpha} & -\widehat{\Sigma} \\
    \Sigma & \alpha \end{pmatrix}.
\]

\begin{lemma}
    \label{lem:kani}
    $\widehat{F} F = F \widehat{F} = 2^e$. Hence $F$ is an isogeny of degree
    $2^{4e}$, and $\ker F = \widehat{F}\bigl(A[2^e]\bigr)$.
\end{lemma}

\begin{proof}
    Every entry of $F$ and $\widehat{F}$ is an integer matrix or a
    componentwise map $\operatorname{diag}(\psi, \psi)$, and these commute with
    each other: for an integer matrix $M$ and any homomorphism $\psi$ of
    curves, $\operatorname{diag}(\psi, \psi) M = M \operatorname{diag}(\psi,
    \psi)$, because $\psi$ is a group homomorphism. The diagonal blocks of
    $\widehat{F} F$ are $\widehat{\alpha}\alpha + \widehat{\Sigma}\Sigma =
    (2^e - q) + q = 2^e$ and $\Sigma\widehat{\Sigma} + \alpha\widehat{\alpha}
    = 2^e$, by Lemma~\ref{lem:J}(iii) and $\widehat{\sigma}\sigma =
    \sigma\widehat{\sigma} = q$; the off-diagonal blocks are
    $\widehat{\alpha}\widehat{\Sigma} - \widehat{\Sigma}\widehat{\alpha} = 0$
    and $\Sigma\alpha - \alpha\Sigma = 0$. The same holds for
    $F\widehat{F}$. Then $(\deg F)^2 = \deg [2^e] = 2^{8e}$ on a fourfold, so
    $\deg F = 2^{4e}$. Finally $\widehat{F}(A[2^e]) \subseteq \ker F$ because
    $F\widehat{F} = 2^e$, and both groups have $2^{4e}$ elements, since
    $\ker \widehat{F} \subseteq A[2^e]$ has $2^{4e}$ elements and
    $|A[2^e]| = 2^{8e}$.
\end{proof}

\begin{proposition}
    \label{prop:commute}
    $F \circ D = D \circ F$. In particular $\ker F$ is stable under $D$.
\end{proposition}

\begin{proof}
    Since $D = \operatorname{diag}(J, J)$, the identity $FD = DF$ holds block
    by block exactly when each block of $F$ commutes with $J$. The blocks
    $\alpha$ and $\widehat{\alpha}$ are polynomials in $J$
    (Lemma~\ref{lem:J}(ii)), and $\Sigma$, $\widehat{\Sigma}$ are componentwise
    maps, which commute with every integer matrix (proof of
    Lemma~\ref{lem:kani}). Explicitly, $J\Sigma(P, Q) = J(\sigma P, \sigma Q)
    = (\sigma Q, -\sigma P) = \Sigma(Q, -P) = \Sigma J(P, Q)$. If $F(x) = 0$
    then $F(D x) = D F(x) = 0$.
\end{proof}

\subsection{Consequences along the chain}

The verifier computes $F$ as a chain of $e$ isogenies of degree 2. Write
$K = \ker F$ and, for $0 \le k \le e$, let $F_k \colon A \to A_k =
A / [2^{e-k}]K$ be the partial chain after $k$ steps, so that $F_e = F$ up
to isomorphism. The two half-chains that meet in the middle are partial
chains of this kind, one read from each end.

\begin{proposition}
    \label{prop:descent}
    For every $k$, there is an automorphism $D_k$ of $A_k$ with
    $D_k \circ F_k = F_k \circ D$. It satisfies $D_k^2 = -1$ and
    $\widehat{D_k} = -D_k$, and it preserves the principal polarisation of
    $A_k$.
\end{proposition}

\begin{proof}
    $D$ is a group automorphism with $D(K) = K$
    (Proposition~\ref{prop:commute}), hence $D([2^{e-k}]K) = [2^{e-k}]K$, so
    $F_k \circ D$ vanishes on $\ker F_k$ and factors as $D_k \circ F_k$; $D_k$
    is an automorphism because $D$ is. The identities $D_k^2 = -1$ and
    $\widehat{D_k} = -D_k$ follow from $D^2 = -1$ and $\widehat{D} = -D$
    (Lemma~\ref{lem:J}) by uniqueness of the factorisation. The principal
    polarisation of $A_k$ is the one induced from $2^k$ times that of $A$
    through the isotropic kernel $\ker F_k$; $D$ preserves both, so $D_k$
    preserves the induced one.
\end{proof}

\begin{remark}[The module viewpoint]
\label{rem:module}
The varieties here are maximal: each curve has $(p+1)^2$ points over
$\mathbb{F}_{p^2}$. By the anti-equivalence $A \mapsto M_A =
\operatorname{Hom}(A, E_0)$ of~\cite[Theorem 2.5]{PRS26}, under which
principal polarisations correspond to unimodular Hermitian forms $H_A$ and
isogenies to module monomorphisms~\cite[Table 1]{PRS26}, the automorphism
$D$ corresponds to $f \mapsto f \circ D$ on $M_A$, which preserves $H_A$
by Lemma~\ref{lem:J}(iv) and squares to $-1$. The partial chain $F_k$
corresponds to the inclusion $\psi_k \colon M_{A_k} \hookrightarrow M_A$,
$f \mapsto f \circ F_k$, whose image consists of the morphisms vanishing on
$\ker F_k$. Since $D(\ker F_k) = \ker F_k$, that image is stable under
$f \mapsto f \circ D$, and $f \circ D_k \circ F_k = f \circ F_k \circ D$
shows that $D_k$ corresponds to the restriction of $f \mapsto f \circ D$
to it. This is Proposition~\ref{prop:descent} in module form: $D$
descends along the chain exactly because each partial kernel, equivalently
each image submodule, is $D$-stable.
\end{remark}

\begin{proposition}
    \label{prop:weil}
    Every $A_k$ is of Weil type in the sense of~\cite[Definition 4.1]{BDD+26}:
    $D_k$ embeds $\mathbb{Z}[i]$ in $\operatorname{End}(A_k)$, with
    $i \mapsto D_k$, and acts on the tangent space at the origin with
    eigenvalues $i, i, -i, -i$.
\end{proposition}

\begin{proof}
    On the tangent space of $E^2$ the differential of $J$ is the matrix
    $\left(\begin{smallmatrix} 0 & 1 \\ -1 & 0 \end{smallmatrix}\right)$,
    whose eigenvalues are $i$ and $-i$; these lie in $\mathbb{F}_{p^2}$ because
    $p \equiv 3 \pmod 4$. Hence $D = \operatorname{diag}(J, J)$ has eigenvalues
    $i, i, -i, -i$ on the tangent space of $A$. Isogenies of degree a power of
    2 are separable in odd characteristic, so the differential of $F_k$ is an
    isomorphism, and it intertwines $D$ with $D_k$
    (Proposition~\ref{prop:descent}); the eigenvalues are therefore the same.
    This is the argument of~\cite[proof of Proposition G.1]{BDD+26}.
\end{proof}

\begin{proposition}
    \label{prop:forged}
    For any matrix the verifier receives, honest or not, the kernel it builds
    is stable under $D$. Whenever that kernel defines a chain, that is, when it
    is maximal isotropic and the computation succeeds, the conclusions of
    Propositions~\ref{prop:descent} and~\ref{prop:weil} hold for that chain.
\end{proposition}

\begin{proof}
    The verifier builds its kernel from the transmitted action rather than
    from $\sigma$: the transmitted matrix defines a map $\sigma'$ on the
    $2^r$-torsion, and the verifier uses $\Sigma' = \operatorname{diag}(\sigma',
    \sigma')$ in place of $\Sigma$. Like $\Sigma$, the map $\Sigma'$ acts
    componentwise, so it commutes with $J$, and the kernel generators
    $(\widehat{\alpha}(x), \Sigma'(x))$ satisfy $D(\widehat{\alpha}(x),
    \Sigma'(x)) = (\widehat{\alpha}(Jx), \Sigma'(Jx))$, which is again of that
    form. The proofs of Propositions~\ref{prop:descent} and~\ref{prop:weil} use
    only that the kernel is $D$-stable and that the chain has 2-power degree.
\end{proof}

\subsection{The kernel as a module over \texorpdfstring{$\mathbb{Z}[i]$}{Z[i]}}

Let $R_e = (\mathbb{Z}/2^e)[i]$, acting on $2^e$-torsion with $i$ acting as
$J$ (or $D$).

\begin{lemma}
    \label{lem:free}
    $E^2[2^e]$ is a free $R_e$-module of rank 2, and so is $\ker F$.
\end{lemma}

\begin{proof}
    Let $(S, T)$ be a basis of $E[2^e]$. Since $J(S, 0) = (0, -S)$ and
    $J(T, 0) = (0, -T)$, the four points
    $(S, 0), (T, 0), (0, S), (0, T)$, which form a $\mathbb{Z}/2^e$-basis of
    $E^2[2^e]$, are $(S, 0)$, $(T, 0)$, $-J(S, 0)$ and $-J(T, 0)$; so $(S, 0)$ and $(T, 0)$ form an $R_e$-basis. For the kernel,
    consider
    \[
    \iota \colon A_1[2^e] \to A, \qquad
    \iota(x) = \widehat{F}(x, 0) = \bigl(\widehat{\alpha}(x), \Sigma(x)\bigr).
    \]
    Its image lies in $\ker F$ by
    Lemma~\ref{lem:kani}. It is injective, because $\widehat{\alpha}$ has odd
    degree $(2^e - q)^2$ and is therefore injective on 2-power torsion. Its
    domain has $2^{4e}$ elements, as does $\ker F$, so $\iota$ is a bijection
    onto $\ker F$. Finally $\iota$ commutes with $J$, so it is $R_e$-linear, and
    $\ker F \cong A_1[2^e] = \Echl^2[2^e]$ is free of rank 2.
\end{proof}

MIKE uses the corresponding statement about its own kernel~\cite[proof of
Corollary G.1]{BDD+26} to show that the pairing $e_{1+i}$ is alternate on
its fourfolds, which in turn is what makes its adapted theta coordinates
exist. On the starting fourfold $A$ the analogous pairing is alternate.

\begin{lemma}
    \label{lem:alternate}
    The Weil--Cartier pairing $e_{1+D}$ on $A[1 + D]$ is alternate.
\end{lemma}

\begin{proof}
    It suffices to treat one square $E^2$, since $A$ is a product of two
    and the pairings multiply. If $(1 + J)(P, Q) = (P + Q, Q - P) = 0$, then
    $P = -Q$ and $2P = 0$, so $\ker(1 + J) = \{(u, u) : u \in E[2]\}$. With
    the formula of~\cite[proof of Corollary G.1]{BDD+26},
    $e_{1+J}(x, x) = e_2(J x_0, x_0)$ for any $x_0$ with $(1 + J) x_0 = x$.
    For $x = (u, u)$ take $x_0 = (0, u)$, since $(1 + J)(0, u) = (0, u) +
    (u, 0) = (u, u)$. Then $e_2(J x_0, x_0) = e_2\bigl((u, 0), (0, u)\bigr) =
    e_2(u, 0)\, e_2(0, u) = 1$.
\end{proof}

\subsection{What remains open}

MIKE's saving of two kernel points out of four~\cite[Theorem
E.1]{BDD+26} does not follow from Propositions~\ref{prop:descent}
and~\ref{prop:weil} alone. It needs level-2 theta coordinates that are
\emph{adapted to $[i]$} in the sense of~\cite[Definition G.1]{BDD+26},
built from a basis of the 4-torsion through the $(1 + i)$-torsion. For
such coordinates, \cite[Proposition G.3]{BDD+26} shows that $[i]$ acts on
the sixteen theta coordinates as a fixed permutation that fixes four of
them and swaps the other twelve in six pairs; that statement assumes only
an adapted basis, so it would apply to $D_k$ unchanged. What MIKE then
proves for its own chain (its Theorem 4.1, Lemma G.3 and gluing
algorithms) is that adapted coordinates can be chosen at the start and
remain adapted at every step. For the verifier of
Section~\ref{sec:dim4}, the following steps are not yet done.
\begin{enumerate}
    \item Choose theta coordinates on $A$, adapted to $D$ and compatible
    with the kernel $\ker F$; Lemmas~\ref{lem:free}
    and~\ref{lem:alternate} are the starting point.
    \item Carry them through the gluing steps, from a product of four
    elliptic curves; these steps differ from MIKE's, which glue into
    Scholten's construction.
    \item Check that the theta coordinates which MIKE assumes nonzero
    (\cite[Heuristic E.1]{BDD+26}) are nonzero on this family of
    fourfolds as well.
\end{enumerate}
For the two-point evaluation of the kernel points, item 1 is settled in
practice by the observation below; items 2 and 3 matter for MIKE's
computation of the codomain from two points, which we have not
implemented.
Because Proposition~\ref{prop:forged} lets a verifier rely on the symmetry
without checking it, a forged signature could still be built to make one
of these coordinates vanish; an optimised verifier must then reject the
signature rather than divide by zero.

\paragraph{What the implementation shows}
We solved for the matrix $N_k$ with $N_k\,\theta(x) \sim \theta(Dx)$ from
20 random points $x$ at the gluing codomain and after plain steps 1, 2, 3,
20, 43, 70 and 84 of the first half-chain of a level-I signature. At every
step the solution is unique up to scale and agrees with held-out points,
as Proposition~\ref{prop:descent} predicts, and it is the same pure
permutation, with all scalars equal to 1: it exchanges index bits 0 and 2
and bits 1 and 3, fixing coordinates 0, 5, 10 and 15 and exchanging the
other twelve in six pairs. This is the shape of~\cite[Proposition
G.3]{BDD+26} up to a relabelling of the index bits, and the theta null
points accordingly have six pairs of equal coordinates. After the gluing,
$D$ maps the first two kernel generators to the other two, up to sign. So
the verifier's coordinates are adapted to $D$ without any change of theta
structure. We have not proved that this holds for every signature
(Conjecture~\ref{conj:twopoint} below); the implementation checks the relation on the kernel generators of each
signature and falls back to four points when it fails, which happened on
none of the 24 signatures tested.

%% ------------------------------------------------------------
%% OPEN ISSUE, set in red on purpose: remove the colour once the
%% conjecture is proved (or the two-point claim is withdrawn).
%% ------------------------------------------------------------
{\color{red}
    \subsection{The conjecture behind the two-point verifier}
    \label{app:conjecture}

    The two-point verifier rests on more than
    Propositions~\ref{prop:commute} to~\ref{prop:forged}: those show that the
    chain respects $D$, not that $D$ acts in the verifier's own theta
    coordinates as one fixed permutation. We state what is needed as a
    conjecture.

    Let $\pi$ be the permutation of $\{0, \ldots, 15\}$ that exchanges bits 0
    and 2 and bits 1 and 3 of the index $k = i_0 + 2 i_1 + 4 i_2 + 8 i_3$.

    \begin{conjecture}
        \label{conj:twopoint}
        For every signature, honest or not, for which the verifier's half-chains
        are defined, and with the symplectic matrices and gluing changes of basis
        the verifier uses:
        \begin{enumerate}[label=(\roman*)]
            \item on every intermediate fourfold $A_k$ of either half-chain,
            including the codomain of the gluing, $\theta_a(D_k x) =
            \theta_{\pi(a)}(x)$ for all $x \in A_k$ and all $a$, as projective
            points in the level-2 theta coordinates the verifier computes;
            \item after the gluing, the four kernel generators satisfy
            $\theta(g_2) = \pi\,\theta(g_0)$ and $\theta(g_3) = \pi\,\theta(g_1)$
            projectively, where $\pi$ permutes coordinates.
        \end{enumerate}
        Consequently, the two-point chain computes the same isogenies, and the
        same codomains, as the four-point chain.
    \end{conjecture}

    \paragraph{Evidence}
    Part (i) was first observed at the gluing codomain and after plain steps
    1, 2, 3, 20, 43, 70 and 84 of the first half-chain of one level-I
    signature, where solving for the matrix of $D_k$ gives $\pi$ and nothing
    else. It was then checked at every step of both half-chains, on four
    random points per signature and on every codomain null point: on 32
    honest level-I signatures from four keys, 64 half-chains and 5\,472
    steps, and on 96 signatures forged by replacing the transmitted response
    entries with random residues, of which 76 pass the setup, giving 144
    computed half-chains and 12\,318 steps. No step failed, and as a control,
    $\theta(Dx)$ never equalled $\theta(x)$ without the permutation. Part
    (ii), and identical codomains along both half-chains, was observed on 24
    level-I signatures from three keys. The forged signatures keep an honestly
    generated $q$, commitment curve and challenge; signatures with a forged
    $q$, and levels III and V, which are not implemented, remain untested.

    \paragraph{Why it matters}
    The implementation checks part (ii) on each signature and falls back to
    four points when it fails, but it does not check part (i) after the
    gluing. If part (i) failed at a later step, the two-point chain would
    build wrong kernels, and the soundness of the verifier, which assumes the
    chain computes the isogeny with the transmitted kernel, would no longer
    follow. Until the conjecture is proved, a cheap safeguard is to check at
    every step that the codomain null point is fixed by $\pi$ (six
    comparisons), a necessary condition that would at least locate a
    failure.

    \paragraph{A tentative approach}
    An induction along the chain, after MIKE's treatment of its own
    fourfolds~\cite[Proposition G.3 and Lemma G.3]{BDD+26}.
    \begin{enumerate}
        \item \emph{The product.} On $E^2$, give both copies the level-2 theta
        structure induced by the same basis of $E[4]$, so that
        $\theta_{(i,j)}(P, Q) = \theta_i(P)\,\theta_j(Q)$. Since level-2 theta
        functions are even,
        \[
        \theta_{(i,j)}\bigl(J(P, Q)\bigr) = \theta_i(Q)\,\theta_j(-P)
        = \theta_{(j,i)}(P, Q),
        \]
        so $J$ exchanges the two index bits of the square, and $D$ on $\Ecom^2
        \times \Echl^2$ exchanges two pairs of index bits. To check: that the
        verifier gives both copies of each curve the same basis, and which
        index bits its product coordinates assign to which copy; $\pi$ suggests
        that the two copies of a curve sit at bits $\{0, 2\}$ and $\{1, 3\}$,
        possibly after the gluing's change of basis.
        \item \emph{The gluing.} The theta structure on the gluing codomain is
        induced through the change-of-basis matrices built from the verifier's
        symplectic matrices. To show: $D$ acts on the resulting symplectic
        basis by a block matrix $\left(\begin{smallmatrix} A & 0 \\ 0 &
            A^{-\mathsf{T}} \end{smallmatrix}\right)$ with $A$ the permutation
        matrix of $\pi$. Such a matrix acts on level-2 theta coordinates by
        permuting their indices, with no characters~\cite[Theorem 12]{Dar24};
        that would give part (i) at the gluing codomain. The Kani blocks are of
        the form $a + bJ$ (Lemma~\ref{lem:J}) and commute with $D$; what needs
        checking is the symplectic completion and the gluing change of basis
        computed from them.
        \item \emph{The induction.} If $D$ acts by that block matrix on the
        symplectic basis inducing the theta structure of $A_k$, then, since
        each step commutes with $D$ (Proposition~\ref{prop:descent}), it acts
        by the same matrix on the image basis inducing the structure of
        $A_{k+1}$, hence by the same permutation $\pi$. The second half-chain,
        computed from the other end, would follow in the same way from its own
        starting matrix.
        \item \emph{The kernel generators.} Part (ii) would follow from the
        form of the kernel basis built from the starting symplectic matrix:
        the kernel of $F$ is $D$-stable for any received matrix
        (Proposition~\ref{prop:forged}), and the question is only whether the
        chosen generators come in pairs $(g, Dg)$.
    \end{enumerate}
    The base cases, steps 1 and 2, depend on the exact matrices in the
    implementation, and are where the work lies.
}
